%═══════════════════════════════════════════════════════════════════════════════
\chapter{Cálculo de Itô: integración estocástica y ecuaciones diferenciales
estocásticas}
\label{ch:cap10}
%═══════════════════════════════════════════════════════════════════════════════

El Capítulo~\ref{ch:cap5} introdujo el cálculo estocástico con urgencia pragmática: la
tabla $(dW_t)^2 = dt$, la fórmula de Itô enunciada, la SDE de Langevin del
SGD usada para derivar la distribución de Gibbs. Fue suficiente para el
propósito de ese capítulo. Pero dejó deudas técnicas precisas que este capítulo
salda.

¿Por qué $(dW_t)^2 = dt$ y no cero? Porque el movimiento browniano tiene
variación cuadrática no trivial, y ese es un resultado que se puede demostrar
con rigor. ¿Por qué la fórmula de Itô tiene el término de corrección de segundo
orden $\frac{\sigma^2}{2}\partial_{xx}f$ que no existe en el cálculo ordinario?
Porque la regla de Taylor estocástica retiene el término de segundo orden en el
incremento $dW$, que contribuye al primer orden en $dt$ precisamente por la
variación cuadrática. ¿Bajo qué condiciones tiene solución una SDE? Bajo
condiciones de Lipschitz y crecimiento lineal, por el argumento de punto fijo
de Picard adaptado al caso estocástico.

La conexión con la física matemática del estudiante es más profunda de lo que
aparenta. El generador infinitesimal de un proceso de difusión es el operador
diferencial de la ecuación de Fokker-Planck; para difusión pura sin deriva, ese
operador es el laplaciano de la ecuación de calor. Las soluciones de las EDPs
parabólicas son densidades de transición de procesos estocásticos, y viceversa.
El cálculo de Itô y las EDPs parabólicas son la misma matemática en dos idiomas
distintos.

%───────────────────────────────────────────────────────────────────────────────
\section{El movimiento browniano: construcción y propiedades}
\label{sec:browniano}
\dificultad{2}{2}
%───────────────────────────────────────────────────────────────────────────────

\subsection{Límite de caminatas aleatorias}
\label{subsec:caminata}

Sea $\xi_1, \xi_2, \ldots$ una sucesión de variables aleatorias i.i.d.\ con
$\mathbb{E}[\xi_k] = 0$ y $\mathrm{Var}[\xi_k] = 1$. La caminata aleatoria
es $S_n = \sum_{k=1}^n \xi_k$. Por el teorema central del límite, $S_n/\sqrt{n}
\to \mathcal{N}(0,1)$ en distribución. El proceso en tiempo continuo
$W^{(n)}(t) = S_{\lfloor nt \rfloor}/\sqrt{n}$ interpola la caminata a
escala temporal $1/n$ y espacial $1/\sqrt{n}$.

El \textbf{teorema de Donsker} (invarianza del principio de invarianza,
enunciado sin demostración) establece que $W^{(n)} \to W$ en el espacio
$C([0,T])$ de funciones continuas con la métrica del supremo, donde $W$ es
el movimiento browniano estándar. La convergencia es en distribución como
proceso: no solo las distribuciones marginales convergen, sino la ley completa
del proceso.

Para el SGD: los incrementos del gradiente estocástico en cada mini-lote son
sumas de términos i.i.d.\ (las contribuciones de las muestras del lote). En el
régimen de muchas iteraciones con lotes pequeños, el proceso acumulado de
iteraciones converge al movimiento browniano escalado. La SDE de Langevin del
Capítulo~\ref{ch:cap5} es exactamente ese límite.

\subsection{Definición axiomática del movimiento browniano}
\label{subsec:browniano_def}

\begin{definicion}{Movimiento browniano estándar}
\label{def:browniano}
Un proceso estocástico $W: [0,\infty) \times \Omega \to \mathbb{R}$ definido
sobre un espacio de probabilidad $(\Omega, \mathcal{F}, \mathbb{P})$ es un
\textbf{movimiento browniano estándar} si satisface:
\begin{enumerate}
  \item $W(0) = 0$ casi seguramente (c.s.).
  \item \textbf{Incrementos independientes:} para $0 \leq t_0 < t_1 < \cdots
        < t_n$, las variables $W(t_1)-W(t_0), W(t_2)-W(t_1),\ldots,
        W(t_n)-W(t_{n-1})$ son independientes.
  \item \textbf{Incrementos estacionarios gaussianos:}
        $W(t)-W(s) \sim \mathcal{N}(0,t-s)$ para $0 \leq s < t$.
  \item \textbf{Continuidad de trayectorias:} las trayectorias $t \mapsto W(t,\omega)$
        son casi seguramente continuas.
\end{enumerate}
\end{definicion}

Las propiedades 1--3 determinan completamente la distribución finito-dimensional
del proceso. La existencia de una versión con trayectorias continuas (propiedad 4)
es el contenido no trivial del \textbf{teorema de Kolmogorov-Chentsov}: si un
proceso satisface $\mathbb{E}[|W(t)-W(s)|^\alpha] \leq C|t-s|^{1+\beta}$ para
constantes $\alpha, \beta, C > 0$, entonces admite una versión con trayectorias
Hölder-continuas de exponente $\gamma < \beta/\alpha$. Para el movimiento
browniano: $\mathbb{E}[|W(t)-W(s)|^4] = 3(t-s)^2$, con $\alpha = 4$ y $\beta = 1$,
lo que garantiza trayectorias Hölder con exponente $\gamma < 1/4$. (De hecho,
las trayectorias son Hölder con exponente $\gamma < 1/2$, pero esto requiere
un argumento más cuidadoso.)

\subsection{No diferenciabilidad y variación cuadrática}
\label{subsec:variacion_cuadratica}

\begin{proposicion}{No diferenciabilidad del movimiento browniano}
\label{prop:no_diff}
Las trayectorias del movimiento browniano no son diferenciables en ningún punto,
casi seguramente.
\end{proposicion}

\begin{proof}
Supóngase que $W(\cdot,\omega)$ es diferenciable en $t$. Entonces $W'(t,\omega) =
\lim_{h\to 0}(W(t+h,\omega)-W(t,\omega))/h$ existe y es finito. Para cualquier
$h > 0$ fijo, $(W(t+h)-W(t))/h \sim \mathcal{N}(0, 1/h)$: tiene varianza $1/h$,
que diverge cuando $h \to 0$. Ninguna sucesión de variables con varianza
divergente puede converger casi seguramente a un límite finito. Más precisamente,
usando el criterio de Cauchy: para que el límite exista es necesario que para
cualquier $\varepsilon > 0$, $\mathbb{P}(|(W(t+h)-W(t))/h - (W(t+h')-W(t))/h'|
\geq \varepsilon) \to 0$ para $h, h' \to 0$. El cálculo de esta probabilidad
usando la distribución gaussiana da un valor que no tiende a cero uniformemente
en $t$: para casi todo $\omega$, hay puntos $t$ donde la diferencia permanece
grande. $\square$
\end{proof}

La consecuencia es que el cálculo ordinario, que requiere diferenciabilidad de
primer orden, no puede aplicarse al movimiento browniano. La noción que lo
reemplaza es la variación cuadrática.

\begin{definicion}{Variación cuadrática}
\label{def:vc}
Para una partición $\Pi = \{0 = t_0 < t_1 < \cdots < t_n = T\}$ con malla
$\|\Pi\| = \max_i(t_{i+1}-t_i)$, la \textbf{variación cuadrática} del
movimiento browniano en $[0,T]$ es el límite de:
\[
  V_T^{(2)}(\Pi) = \sum_{i=0}^{n-1}(W(t_{i+1})-W(t_i))^2.
\]
\end{definicion}

\begin{teorema}{Variación cuadrática del movimiento browniano}
\label{thm:vc}
Para cualquier sucesión de particiones $\{\Pi_n\}$ con $\|\Pi_n\| \to 0$:
\[
  V_T^{(2)}(\Pi_n) \xrightarrow{c.s.} T.
\]
\end{teorema}

\begin{proof}
Sean $\Delta W_i = W(t_{i+1})-W(t_i)$ y $\Delta t_i = t_{i+1}-t_i$.
Como $\Delta W_i \sim \mathcal{N}(0,\Delta t_i)$, se tienen:
\[
  \mathbb{E}[(\Delta W_i)^2] = \Delta t_i,
  \qquad
  \mathrm{Var}[(\Delta W_i)^2] = 2(\Delta t_i)^2.
\]
La suma es $V_T^{(2)} = \sum_i (\Delta W_i)^2$, con:
\[
  \mathbb{E}[V_T^{(2)}] = \sum_i \Delta t_i = T,
\]
\[
  \mathrm{Var}[V_T^{(2)}]
  = \sum_i \mathrm{Var}[(\Delta W_i)^2]
  = 2\sum_i(\Delta t_i)^2
  \leq 2\|\Pi\|\sum_i \Delta t_i = 2\|\Pi\|T \to 0.
\]
Esto establece convergencia en $L^2$: $V_T^{(2)} \to T$ en $L^2(\mathbb{P})$.

Para la convergencia \textbf{casi segura}, se aplica el segundo lema de
Borel-Cantelli. Sea $Z_n = V_T^{(2)}(\Pi_n) - T$. De la estimación anterior,
$\mathrm{Var}[Z_n] \leq 2\|\Pi_n\|T$. Para particiones diádicas con $n$
subintervalos de longitud $\Delta t = T/n$:
$\mathrm{Var}[Z_n] = 2T^2/n$.
Por la desigualdad de Markov y Chebyshev:
\[
  \mathbb{P}(|Z_n| \geq \varepsilon)
  \leq \frac{\mathrm{Var}[Z_n]}{\varepsilon^2}
  = \frac{2T^2}{n\varepsilon^2}.
\]
Esta cota es summable: $\sum_{n=1}^\infty \mathbb{P}(|Z_n| \geq \varepsilon) \leq
\frac{2T^2}{\varepsilon^2}\sum_{n=1}^\infty \frac{1}{n} = \infty$: no sumable.
Para obtener convergencia c.s.\ es necesario acotar más finamente.

Se toma en cambio la subsucesión $n_k = k^2$ (partición en $k^2$ intervalos):
\[
  \mathbb{P}(|Z_{n_k}| \geq \varepsilon)
  \leq \frac{2T^2}{k^2\varepsilon^2},
\]
cuya suma $\sum_{k=1}^\infty \frac{1}{k^2} < \infty$ es convergente. Por el
segundo lema de Borel-Cantelli:
$\mathbb{P}(|Z_{n_k}| \geq \varepsilon \text{ i.o.}) = 0$,
es decir, $Z_{n_k} \to 0$ c.s.\ a lo largo de la subsucesión $\{n_k\}$.

Para pasar de la subsucesión a la sucesión completa: para $n_k \leq n \leq
n_{k+1}$, los términos adicionales en $V^{(2)}(\Pi_n) - V^{(2)}(\Pi_{n_k})$
se controlan acotando el oscilación máxima del browniano en subintervalos
de tamaño $O(1/n)$. Usando la continuidad uniforme de las trayectorias
(el browniano es uniformemente continuo en $[0,T]$ c.s.) y el hecho de que
hay $O(k^2)$ subintervalos adicionales de tamaño $O(1/k^2)$, la diferencia
es $O(1/k)$ c.s., lo que completa el argumento. $\square$
\end{proof}

\begin{observacion}{La notación $(dW_t)^2 = dt$}
El Teorema~\ref{thm:vc} es la justificación precisa de la regla mnemotécnica
$(dW_t)^2 = dt$ usada en el Capítulo~\ref{ch:cap5}. No es una regla de álgebra: es el
enunciado de que la variación cuadrática del movimiento browniano en $[0,t]$
es exactamente $t$. En el lenguaje diferencial, incrementar el browniano en
$dW_t$ produce una contribución $dW_t \cdot dW_t = dt$ al segundo orden que,
a diferencia del cálculo ordinario, no es $o(dt)$ sino exactamente $dt$.
Esta es la razón fundamental por la que el cálculo estocástico difiere del
cálculo ordinario.
\end{observacion}

\begin{fisicaconexion}{Movimiento browniano y difusión en física}
El movimiento browniano como objeto físico fue observado por Robert Brown en
1827 (polen en agua) y explicado teóricamente por Einstein en 1905 y Smoluchowski
en 1906. Einstein derivó que el desplazamiento cuadrático medio de una partícula
difusiva satisface $\langle x^2(t)\rangle = 2Dt$, donde $D = k_BT/(6\pi\eta r)$
es el coeficiente de difusión (relación de Einstein-Stokes): $k_B$ la constante
de Boltzmann, $T$ la temperatura, $\eta$ la viscosidad y $r$ el radio de la
partícula. Esto es exactamente el Teorema~\ref{thm:vc}: la variación cuadrática
$\langle(W(t))^2\rangle = t$ con el coeficiente de difusión absorbido en la
escala de $W$.

La construcción del Teorema~\ref{thm:vc} tiene una consecuencia directa para
la medición experimental. La variación cuadrática $V_T^{(2)}(\Pi) = \sum_i(\Delta
x_i)^2$ es un estimador consistente del coeficiente de difusión: observar una
trayectoria discreta con resolución temporal $\Delta t$ permite estimar $D$
por $\hat{D} = V_T^{(2)}/(2T)$. En espectroscopía de correlación de fluorescencia
(FCS) y seguimiento de partícula única (SPT), este principio se usa para medir
la difusión de proteínas en membranas celulares a escala nanométrica. La
Proposición~\ref{prop:no_diff} —no diferenciabilidad casi segura— implica que
la \emph{velocidad instantánea} de la partícula browniana no existe: solo
tiene sentido el desplazamiento cuadrático medio, no la velocidad. Este es el
límite de sobreenfriamiento viscoso de la ecuación de Langevin completa
$m\ddot{x} = -\gamma\dot{x} + F_{\mathrm{ruido}}$: cuando $m/\gamma \to 0$,
la inercia desaparece y la velocidad se vuelve singular.
\end{fisicaconexion}
\label{subsec:martingala}

\begin{definicion}{Filtración y adaptabilidad}
\label{def:filtracion}
La \textbf{filtración natural} del movimiento browniano es la familia creciente
de $\sigma$-álgebras $\{\mathcal{F}_t\}_{t \geq 0}$ con
$\mathcal{F}_t = \sigma(W(s): 0 \leq s \leq t)$: la información acumulada
hasta el tiempo $t$. Un proceso estocástico $X_t$ es \textbf{adaptado} a
$\{\mathcal{F}_t\}$ si $X_t$ es $\mathcal{F}_t$-medible para todo $t$: el
valor de $X_t$ es determinado por la trayectoria de $W$ hasta $t$, sin
información del futuro.
\end{definicion}

\begin{definicion}{Martingala}
\label{def:martingala}
Un proceso adaptado $M_t$ con $\mathbb{E}[|M_t|] < \infty$ es una
\textbf{martingala} respecto a $\{\mathcal{F}_t\}$ si:
\[
  \mathbb{E}[M_t | \mathcal{F}_s] = M_s \quad \text{para todo } s \leq t.
\]
\end{definicion}

Tres martingalas fundamentales del movimiento browniano (verificación en
el Ejercicio~\ref{ej:martingalas}):
\begin{align}
  M_t^{(1)} &= W_t, \label{eq:mart1}\\
  M_t^{(2)} &= W_t^2 - t, \label{eq:mart2}\\
  M_t^{(3)} &= e^{\theta W_t - \theta^2 t/2}
  \quad (\text{para cualquier } \theta \in \mathbb{R}).
  \label{eq:mart3}
\end{align}
$M_t^{(1)}$ dice que el mejor predictor del valor futuro del browniano es su
valor presente. $M_t^{(2)}$ dice que el cuadrado del browniano crece en
promedio exactamente como $t$: el drift del cuadrado es la variación cuadrática.
$M_t^{(3)}$ (la exponencial de Wick) es la función generatriz de momentos del
browniano y será el bloque con el que se construye el cambio de medida en la
derivación del proceso inverso de la difusión.

%───────────────────────────────────────────────────────────────────────────────
\section{La integral de Itô}
\label{sec:integral_ito}
\dificultad{3}{2}
%───────────────────────────────────────────────────────────────────────────────

\subsection{Por qué la integral de Riemann-Stieltjes falla}
\label{subsec:fallo_rs}

La integral de Riemann-Stieltjes $\int_0^T f(t)\,dg(t)$ existe cuando
$g$ tiene variación total finita $\sum_i |g(t_{i+1})-g(t_i)| < \infty$.

\begin{proposicion}{El browniano tiene variación total infinita c.s.}
\label{prop:var_total}
Para cualquier $T > 0$, la variación total del movimiento browniano en $[0,T]$
satisface $\sum_i|W(t_{i+1})-W(t_i)| \to \infty$ c.s.\ cuando $\|\Pi\| \to 0$.
\end{proposicion}

\begin{proof}
Sea $V_T^{(1)}(\Pi) = \sum_i |\Delta W_i|$ la variación total y $V_T^{(2)}(\Pi)
= \sum_i (\Delta W_i)^2$ la variación cuadrática. Se tiene:
\[
  V_T^{(2)}(\Pi)
  = \sum_i (\Delta W_i)^2
  \leq \max_i|\Delta W_i| \cdot V_T^{(1)}(\Pi).
\]
Por la continuidad uniforme de las trayectorias, $\max_i|\Delta W_i| \to 0$ c.s.\
cuando $\|\Pi\| \to 0$. Como $V_T^{(2)}(\Pi) \to T > 0$ c.s., se necesita
$V_T^{(1)}(\Pi) \to \infty$ c.s.\ para compensar el factor que tiende a cero.
$\square$
\end{proof}

La integral de Riemann-Stieltjes no puede definirse para el movimiento browniano.
Necesitamos una nueva definición de integral que maneje integradores con variación
cuadrática no trivial.

\subsection{Construcción de la integral de Itô}
\label{subsec:construccion}

La integral de Itô $\int_0^T f(t)\,dW_t$ se construye en tres pasos.

\medskip
\noindent\textbf{Paso 1: procesos simples (escalonados).}

Un proceso $f$ es \textbf{simple} si $f(t) = \sum_{i=0}^{n-1} f_i
\mathbf{1}_{[t_i,t_{i+1})}(t)$ donde $f_i$ es $\mathcal{F}_{t_i}$-medible
(adaptado, usando información hasta el inicio del intervalo). Para $f$ simple:
\begin{equation}
  \int_0^T f(t)\,dW_t = \sum_{i=0}^{n-1} f_i(W(t_{i+1})-W(t_i)).
  \label{eq:ito_simple}
\end{equation}
El requisito de que $f_i$ sea $\mathcal{F}_{t_i}$-medible (no $\mathcal{F}_{t_{i+1}}$)
es la diferencia crucial respecto a la integral de Stratonovich y es la fuente
del término de corrección en la fórmula de Itô.

\medskip
\noindent\textbf{Paso 2: la isometría de Itô para procesos simples.}

\begin{teorema}{Isometría de Itô}
\label{thm:isometria}
Para $f$ simple y adaptado:
\begin{equation}
  \mathbb{E}\!\left[\left(\int_0^T f(t)\,dW_t\right)^2\right]
  = \int_0^T \mathbb{E}[f(t)^2]\,dt.
  \label{eq:isometria}
\end{equation}
\end{teorema}

\begin{proof}
Expandiendo el cuadrado de~\eqref{eq:ito_simple}:
\[
  \left(\sum_i f_i\Delta W_i\right)^2
  = \sum_i f_i^2 (\Delta W_i)^2
  + 2\sum_{i < j} f_i f_j \Delta W_i \Delta W_j.
\]
Tomando la esperanza. Para los términos diagonales: $\mathbb{E}[f_i^2(\Delta W_i)^2]
= \mathbb{E}[f_i^2]\,\Delta t_i$ porque $f_i$ es $\mathcal{F}_{t_i}$-medible
e independiente de $\Delta W_i \sim \mathcal{N}(0,\Delta t_i)$.

Para los términos cruzados con $i < j$: $f_j$ es $\mathcal{F}_{t_j}$-medible,
$f_i$ es $\mathcal{F}_{t_i}$-medible, y $\Delta W_j = W(t_{j+1})-W(t_j)$ es
independiente de $\mathcal{F}_{t_j}$ (incremento futuro). Por tanto:
$\mathbb{E}[f_i f_j \Delta W_i \Delta W_j] = \mathbb{E}[f_i f_j \Delta W_i]
\cdot \mathbb{E}[\Delta W_j] = 0$ (ya que $\mathbb{E}[\Delta W_j] = 0$).

Sumando: $\mathbb{E}[(\int_0^T f\,dW)^2] = \sum_i \mathbb{E}[f_i^2]\,\Delta t_i
= \int_0^T \mathbb{E}[f(t)^2]\,dt$. $\square$
\end{proof}

La isometría de Itô establece que la integral es una isometría entre el espacio
$L^2(\Omega \times [0,T])$ de procesos adaptados de cuadrado integrable y
$L^2(\Omega)$. Esto permite extender la integral por completitud.

\medskip
\noindent\textbf{Paso 3: extensión por densidad.}

Sea $\mathcal{L}^2 = \{f: [0,T]\times\Omega \to \mathbb{R}$ adaptado$\,:\,
\int_0^T \mathbb{E}[f(t)^2]\,dt < \infty\}$.
Los procesos simples son densos en $\mathcal{L}^2$. Para cualquier $f \in
\mathcal{L}^2$, existe una sucesión de procesos simples $f^{(n)} \to f$ en
$\mathcal{L}^2$. Por la isometría, $\int_0^T f^{(n)}\,dW$ es una sucesión de
Cauchy en $L^2(\Omega)$:
\[
  \mathbb{E}\!\left[\left(\int_0^T (f^{(n)}-f^{(m)})\,dW\right)^2\right]
  = \int_0^T \mathbb{E}[(f^{(n)}_t-f^{(m)}_t)^2]\,dt \to 0.
\]
El límite en $L^2(\Omega)$ es la \textbf{integral de Itô} $\int_0^T f\,dW$.
La construcción es independiente de la elección de aproximación por procesos
simples.

\subsection{Propiedades fundamentales}
\label{subsec:propiedades_ito}

\begin{proposicion}{Propiedades de la integral de Itô}
\label{prop:ito_propiedades}
Sea $M_t = \int_0^t f(s)\,dW_s$ con $f \in \mathcal{L}^2$. Entonces:
\begin{enumerate}
  \item \textbf{Linealidad:} $\int_0^T(\alpha f+\beta g)\,dW
        = \alpha\int_0^T f\,dW + \beta\int_0^T g\,dW$.
  \item \textbf{Martingala:} $M_t$ es una martingala con $\mathbb{E}[M_t] = 0$.
  \item \textbf{Isometría:} $\mathbb{E}[M_T^2] = \int_0^T\mathbb{E}[f(t)^2]\,dt$.
  \item \textbf{Variación cuadrática:} $[M,M]_t = \int_0^t f(s)^2\,ds$.
\end{enumerate}
\end{proposicion}

\begin{proof}
(2) Para $s < t$, usando la descomposición $M_t = M_s + \int_s^t f(u)\,dW_u$:
\[
  \mathbb{E}[M_t|\mathcal{F}_s]
  = M_s + \mathbb{E}\!\left[\int_s^t f(u)\,dW_u \Big|\mathcal{F}_s\right].
\]
El último término es cero: para el proceso simple, $\mathbb{E}[f_j\Delta W_j|
\mathcal{F}_s] = f_j\mathbb{E}[\Delta W_j|\mathcal{F}_{t_j}] = 0$ porque
$\Delta W_j$ es independiente de $\mathcal{F}_{t_j}$ para $t_j \geq s$. Por
continuidad en $L^2$, el resultado se extiende a $f \in \mathcal{L}^2$.

(4) Para procesos simples: $[M,M]_T = \lim_{\|\Pi\|\to 0} \sum_i(M_{t_{i+1}}
-M_{t_i})^2 = \lim_{\|\Pi\|\to 0}\sum_i f_i^2(\Delta W_i)^2$. Por la
ley de los grandes números para variables independientes con media $f_i^2\Delta t_i$,
la suma converge a $\sum_i f_i^2\Delta t_i = \int_0^T f(t)^2\,dt$. $\square$
\end{proof}

\subsection{La integral de Stratonovich y la diferencia con Itô}
\label{subsec:stratonovich}

La \textbf{integral de Stratonovich} usa el punto medio del intervalo:
\begin{equation}
  \int_0^T f(t) \circ dW_t
  = \lim_{\|\Pi\|\to 0}\sum_i
  \frac{f(t_i)+f(t_{i+1})}{2}(W(t_{i+1})-W(t_i)).
  \label{eq:stratonovich}
\end{equation}
Para $f$ suficientemente regular (por ejemplo, $f = h(W_t)$ para $h \in C^1$),
la diferencia entre las dos integrales es:
\begin{equation}
  \int_0^T f(t) \circ dW_t
  = \int_0^T f(t)\,dW_t + \frac{1}{2}\int_0^T \frac{\partial f}{\partial W}\,dt.
  \label{eq:ito_strat}
\end{equation}
La integral de Stratonovich satisface la regla de la cadena ordinaria $d(g(W_t))
= g'(W_t) \circ dW_t$, sin término de corrección. La integral de Itô satisface
$d(g(W_t)) = g'(W_t)\,dW_t + \frac{1}{2}g''(W_t)\,dt$ (fórmula de Itô).

La elección entre Itô y Stratonovich tiene consecuencias físicas. En sistemas
con ruido multiplicativo (donde el coeficiente de difusión depende del estado),
la interpretación física correcta depende de si el ruido es blanco (límite de
Itô) o de color (límite de Stratonovich). Para los modelos de difusión del
Capítulo~\ref{ch:cap12} se usará la formulación de Itô, que tiene la propiedad de martingala
y es la más natural para el análisis de la distribución estacionaria.

%───────────────────────────────────────────────────────────────────────────────
\section{La fórmula de Itô}
\label{sec:formula_ito}
\dificultad{3}{2}
%───────────────────────────────────────────────────────────────────────────────

\subsection{La regla de la cadena estocástica}
\label{subsec:cadena_estocastica}

En cálculo ordinario, si $X(t)$ es diferenciable y $f$ es de clase $C^1$,
entonces $\frac{d}{dt}f(X(t)) = f'(X(t))X'(t)$, o en notación diferencial:
$df(X) = f'(X)dX$. Esta regla depende de que $(dX)^2 = O((dt)^2)$ sea
despreciable. Para el movimiento browniano, $(dW_t)^2 = dt$ no es despreciable:
el término de segundo orden de la expansión de Taylor contribuye al primer orden.

\subsection{La fórmula de Itô escalar: demostración completa}
\label{subsec:ito_escalar}

\begin{definicion}{Proceso de Itô}
\label{def:ito_proceso}
Un \textbf{proceso de Itô} es un proceso de la forma:
\[
  X_t = X_0 + \int_0^t \mu_s\,ds + \int_0^t \sigma_s\,dW_s,
\]
con $\mu, \sigma \in \mathcal{L}^2$, escrito en forma diferencial como
$dX_t = \mu_t\,dt + \sigma_t\,dW_t$.
\end{definicion}

\begin{teorema}{Fórmula de Itô (caso escalar)}
\label{thm:ito_formula}
Sea $X_t$ un proceso de Itô con $dX_t = \mu_t\,dt + \sigma_t\,dW_t$ y
$f(t,x) \in C^{1,2}([0,T]\times\mathbb{R})$ (una vez diferenciable en $t$,
dos veces en $x$). Entonces $f(t,X_t)$ es también un proceso de Itô y:
\begin{equation}
  df(t,X_t)
  = \left(\frac{\partial f}{\partial t}
    + \mu_t\frac{\partial f}{\partial x}
    + \frac{\sigma_t^2}{2}\frac{\partial^2 f}{\partial x^2}\right)dt
    + \sigma_t\frac{\partial f}{\partial x}\,dW_t.
  \label{eq:ito_formula}
\end{equation}
\end{teorema}

\begin{proof}
Se aplica la expansión de Taylor de orden 2 a $f(t,X_t)$ en el intervalo
$[t_i, t_{i+1}]$ de una partición:
\begin{align*}
  \Delta f_i
  &= f(t_{i+1},X_{t_{i+1}}) - f(t_i,X_{t_i}) \\
  &= \partial_t f\,\Delta t_i
   + \partial_x f\,\Delta X_i
   + \frac{1}{2}\partial_{xx}f\,(\Delta X_i)^2
   + \frac{1}{2}\partial_{tt}f\,(\Delta t_i)^2
   + \partial_{tx}f\,\Delta t_i\,\Delta X_i
   + O(|\Delta|^3),
\end{align*}
donde las derivadas se evalúan en $(t_i, X_{t_i})$ y $\Delta X_i =
X_{t_{i+1}}-X_{t_i} = \mu_{t_i}\Delta t_i + \sigma_{t_i}\Delta W_i$.

Expandiendo $(\Delta X_i)^2$:
\[
  (\Delta X_i)^2
  = \mu_{t_i}^2(\Delta t_i)^2
  + 2\mu_{t_i}\sigma_{t_i}\Delta t_i\,\Delta W_i
  + \sigma_{t_i}^2(\Delta W_i)^2.
\]
Usando la tabla de multiplicación de diferenciales estocásticas:
\begin{center}
\begin{tabular}{ccc}
\toprule
$\times$ & $dt$ & $dW_t$ \\
\midrule
$dt$ & $0$ & $0$ \\
$dW_t$ & $0$ & $dt$ \\
\bottomrule
\end{tabular}
\end{center}
En el límite $\|\Pi\| \to 0$:
\begin{itemize}
  \item $\sum_i (\Delta t_i)^2 \leq \|\Pi\|\sum_i\Delta t_i = \|\Pi\|T \to 0$,
  \item $\sum_i \Delta t_i\,\Delta W_i \to 0$ en $L^2$ (covariación de $t$ y $W$),
  \item $\sum_i (\Delta W_i)^2 \to T$ c.s.\ (Teorema~\ref{thm:vc}).
\end{itemize}

Por tanto, sumando sobre la partición:
\begin{align*}
  f(T,X_T) - f(0,X_0)
  &= \int_0^T \partial_t f\,dt
   + \int_0^T \partial_x f\,\mu_t\,dt
   + \int_0^T \partial_x f\,\sigma_t\,dW_t \\
  &\quad + \frac{1}{2}\int_0^T \partial_{xx}f\,\sigma_t^2\,dt
   + 0 + 0 + O(\|\Pi\|),
\end{align*}
que en el límite da la fórmula~\eqref{eq:ito_formula}. La convergencia es
rigurosa en $L^2$ usando la isometría de Itô para el término estocástico y
la convergencia de sumas de Riemann para los términos deterministas. $\square$
\end{proof}

\begin{observacion}{El término de corrección de Itô}
El término $\frac{\sigma_t^2}{2}\partial_{xx}f$ es la corrección respecto
al cálculo ordinario. Su origen es inequívoco: proviene de $(\Delta W_i)^2
\to dt$, que a su vez proviene de la variación cuadrática no trivial del
movimiento browniano. En el cálculo ordinario, para cualquier función
diferenciable $g(t)$ se tiene $(dg)^2 = O((dt)^2)$ y ese término no contribuye.
Para el browniano, $(dW)^2 = dt$ y contribuye al primer orden. La corrección
de Itô es la marca algebraica de la no diferenciabilidad del browniano.
\end{observacion}

\subsection{La fórmula de Itô multidimensional}
\label{subsec:ito_multi}

Para $\bm{X}_t \in \mathbb{R}^d$ con $d\bm{X}_t = \bm{\mu}_t\,dt +
\bm{\Sigma}_t\,d\bm{W}_t$ ($\bm{W}_t \in \mathbb{R}^m$ browniano $m$-dimensional
con componentes independientes) y $f(t,\bm{x}) \in C^{1,2}$:
\begin{equation}
  df(t,\bm{X}_t)
  = \left(\partial_t f
    + (\nabla_{\bm{x}} f)^\top\bm{\mu}_t
    + \frac{1}{2}\mathrm{tr}\!\left(\bm{\Sigma}_t\bm{\Sigma}_t^\top
    \nabla^2_{\bm{x}} f\right)\right)dt
    + (\nabla_{\bm{x}} f)^\top\bm{\Sigma}_t\,d\bm{W}_t.
  \label{eq:ito_multi}
\end{equation}
El tensor de difusión $\bm{D}_t = \bm{\Sigma}_t\bm{\Sigma}_t^\top \in
\mathbb{R}^{d\times d}$ es semidefinido positivo. El operador
$\frac{1}{2}\mathrm{tr}(\bm{D}\nabla^2 f) = \frac{1}{2}\sum_{i,j} D_{ij}
\partial_{x_i x_j} f$ es el \textbf{generador difusivo}. Para difusión
isótropa $\bm{D} = \sigma^2\bm{I}$: este operador es $\frac{\sigma^2}{2}\Delta f$,
el laplaciano escalado que el estudiante reconoce de la ecuación de calor.

\subsection{Aplicaciones de la fórmula de Itô}
\label{subsec:ito_aplicaciones}

\medskip
\noindent\textbf{Verificación de las martingalas del browniano.}
Aplicando~\eqref{eq:ito_formula} a $f(t,x) = x^2 - t$:
$df = -dt + 2x\,dx + \frac{1}{2}\cdot 2\cdot(dW_t)^2 = -dt + 2W_t\,dW_t + dt
= 2W_t\,dW_t$. El coeficiente de $dt$ es cero: $W_t^2 - t$ es martingala.

Aplicando a $f(t,x) = e^{\theta x - \theta^2 t/2}$:
$df = -\frac{\theta^2}{2}f\,dt + \theta f\,dW_t + \frac{\theta^2}{2}f\,dt
= \theta f\,dW_t$. Coeficiente de $dt$ cero: $e^{\theta W_t - \theta^2 t/2}$
es martingala.

\medskip
\noindent\textbf{La SDE de Langevin y la pérdida del SGD.}
Para $dX_t = -\nabla\mathcal{L}(X_t)\,dt + \sqrt{\eta c}\,dW_t$ y
$f = \mathcal{L}$:
\[
  d\mathcal{L}(X_t) = -\|\nabla\mathcal{L}\|^2\,dt
  + \frac{\eta c}{2}\Delta\mathcal{L}\,dt
  + \sqrt{\eta c}\,\nabla\mathcal{L}\cdot d\bm{W}_t.
\]
En valor esperado: $\frac{d}{dt}\mathbb{E}[\mathcal{L}] = -\mathbb{E}
[\|\nabla\mathcal{L}\|^2] + \frac{\eta c}{2}\mathbb{E}[\Delta\mathcal{L}]$.
El primer término es la disipación (descenso de la pérdida por el gradiente);
el segundo es la inyección de energía por el ruido. En el estado estacionario,
ambos se compensan: $\mathbb{E}[\|\nabla\mathcal{L}\|^2] = \frac{\eta c}{2}
\mathbb{E}[\Delta\mathcal{L}]$, que es la condición de equilibrio de la
distribución de Gibbs del Capítulo~\ref{ch:cap5}.

%───────────────────────────────────────────────────────────────────────────────
\section{Ecuaciones diferenciales estocásticas}
\label{sec:sde}
\dificultad{2}{2}
%───────────────────────────────────────────────────────────────────────────────

\subsection{Definición y soluciones fuertes}
\label{subsec:sde_def}

\begin{definicion}{Ecuación diferencial estocástica}
\label{def:sde}
Una \textbf{SDE} es la ecuación integral:
\begin{equation}
  X_t = X_0 + \int_0^t \mu(X_s,s)\,ds + \int_0^t \sigma(X_s,s)\,dW_s,
  \label{eq:sde}
\end{equation}
con $\mu, \sigma: \mathbb{R}\times[0,T] \to \mathbb{R}$ los coeficientes de
\textbf{deriva} (\textit{drift}) y \textbf{difusión} respectivamente.
Una \textbf{solución fuerte} es un proceso adaptado $X_t$ que satisface~\eqref{eq:sde}
c.s.\ para cada realización del movimiento browniano $W$.
\end{definicion}

La notación diferencial $dX_t = \mu(X_t,t)\,dt + \sigma(X_t,t)\,dW_t$ es una
abreviatura conveniente para la ecuación integral. La SDE en forma diferencial
no tiene significado propio: debe interpretarse como la ecuación integral
donde el segundo término es la integral de Itô.

\subsection{Teorema de existencia y unicidad}
\label{subsec:existencia}

\begin{teorema}{Existencia y unicidad de soluciones fuertes (Itô-Doob)}
\label{thm:ito_doob}
Si los coeficientes satisfacen para constante $L > 0$ y todo $x, y \in \mathbb{R}$,
$t \in [0,T]$:
\begin{enumerate}
  \item \textbf{Lipschitz:}
        $|\mu(x,t)-\mu(y,t)| + |\sigma(x,t)-\sigma(y,t)| \leq L|x-y|$.
  \item \textbf{Crecimiento lineal:}
        $|\mu(x,t)|^2 + |\sigma(x,t)|^2 \leq L^2(1+|x|^2)$.
\end{enumerate}
Entonces para cualquier $X_0$ con $\mathbb{E}[X_0^2] < \infty$, existe una
única solución fuerte de~\eqref{eq:sde} con $\sup_{t\in[0,T]}\mathbb{E}[X_t^2]
< \infty$.
\end{teorema}

\begin{proof}[Esquema (iteración de Picard)]
Se construye la sucesión $X_t^{(0)} = X_0$ y, para $n \geq 0$:
\[
  X_t^{(n+1)} = X_0 + \int_0^t \mu(X_s^{(n)},s)\,ds
  + \int_0^t \sigma(X_s^{(n)},s)\,dW_s.
\]

\noindent\textbf{Convergencia en $L^2$:}
Sea $e_t^{(n)} = \mathbb{E}[\sup_{s \leq t}|X_s^{(n+1)}-X_s^{(n)}|^2]$.
Usando la desigualdad de Doob para martingalas y la isometría de Itô:
\begin{align*}
  e_t^{(n)}
  &\leq C\bigg(\int_0^t \mathbb{E}[|\mu(X_s^{(n)})-\mu(X_s^{(n-1)})|^2]\,ds \\
  &\qquad + \int_0^t \mathbb{E}[|\sigma(X_s^{(n)})-\sigma(X_s^{(n-1)})|^2]\,ds
    \bigg)
  \leq 2CL^2 \int_0^t e_s^{(n-1)}\,ds.
\end{align*}
Por inducción: $e_t^{(n)} \leq \frac{(2CL^2 t)^n}{n!}e_T^{(0)}$.
Como $\sum_n \sqrt{e_T^{(n)}} < \infty$, la sucesión converge en $L^2$
uniformemente en $[0,T]$, y el límite $X_t$ satisface la SDE~\eqref{eq:sde}.

\noindent\textbf{Unicidad:}
Si $X_t$ e $Y_t$ son dos soluciones con la misma condición inicial:
\[
  \mathbb{E}[|X_t-Y_t|^2]
  \leq 2L^2(t+1)\int_0^t \mathbb{E}[|X_s-Y_s|^2]\,ds.
\]
Por el lema de Gronwall: $\mathbb{E}[|X_t-Y_t|^2] \leq 0 \cdot C(t,T) = 0$
para todo $t$, luego $X_t = Y_t$ c.s. $\square$
\end{proof}

\subsection{El proceso de Ornstein-Uhlenbeck: solución explícita}
\label{subsec:ou}

El proceso de Ornstein-Uhlenbeck (OU) es la SDE lineal:
\begin{equation}
  dX_t = -\alpha X_t\,dt + \sigma\,dW_t, \qquad X_0 = x_0,
  \label{eq:ou}
\end{equation}
con $\alpha > 0$ (tasa de retorno) y $\sigma > 0$ (intensidad del ruido).
Los coeficientes $\mu(x) = -\alpha x$ y $\sigma$ cumplen las condiciones del
Teorema~\ref{thm:ito_doob}: la solución existe, es única y puede calcularse
explícitamente.

\medskip
\noindent\textbf{Solución por variación de parámetros (factor integrante).}

\noindent\textit{Paso 1.} El proceso $Y_t = e^{\alpha t}X_t$. Aplicar la
fórmula de Itô al producto $f(t,x) = e^{\alpha t}x$ con $dX_t = -\alpha X_t\,dt
+ \sigma\,dW_t$:
\begin{align*}
  dY_t &= d(e^{\alpha t}X_t) \\
  &= \alpha e^{\alpha t}X_t\,dt + e^{\alpha t}\,dX_t \\
  &= \alpha e^{\alpha t}X_t\,dt + e^{\alpha t}(-\alpha X_t\,dt + \sigma\,dW_t) \\
  &= \sigma e^{\alpha t}\,dW_t.
\end{align*}

\noindent\textit{Paso 2.} Integrando:
$Y_t - Y_0 = \sigma\int_0^t e^{\alpha s}\,dW_s$, es decir:
\begin{equation}
  X_t = x_0 e^{-\alpha t}
  + \sigma\int_0^t e^{-\alpha(t-s)}\,dW_s.
  \label{eq:ou_sol}
\end{equation}

\noindent\textit{Paso 3.} La integral estocástica de un proceso determinista
es una variable gaussiana. Calculando media y varianza:
\begin{align*}
  \mathbb{E}[X_t] &= x_0 e^{-\alpha t}, \\
  \mathrm{Var}[X_t]
  &= \sigma^2\int_0^t e^{-2\alpha(t-s)}\,ds
   = \frac{\sigma^2}{2\alpha}(1-e^{-2\alpha t}).
\end{align*}
Por tanto $X_t \sim \mathcal{N}(x_0 e^{-\alpha t},\,\frac{\sigma^2}{2\alpha}
(1-e^{-2\alpha t}))$.

\medskip
\noindent\textbf{Distribución estacionaria.}
Cuando $t \to \infty$:
\begin{equation}
  X_\infty \sim \mathcal{N}\!\left(0,\, \frac{\sigma^2}{2\alpha}\right).
  \label{eq:ou_estacionaria}
\end{equation}
La distribución estacionaria es gaussiana, centrada en cero, con varianza
$\sigma^2/(2\alpha)$, independiente de la condición inicial $x_0$. La tasa de
convergencia a la estacionaria es $e^{-\alpha t}$: más rápida para $\alpha$
grande (resorte fuerte) y más lenta para $\alpha$ pequeño.

\begin{observacion}{El proceso OU como forward process de la difusión}
El forward process de los modelos de difusión del Capítulo~\ref{ch:cap12} es exactamente
el proceso OU con parámetros específicos. Con $\alpha = 1$ y $\sigma = \sqrt{2}$,
la varianza estacionaria es $\sigma^2/(2\alpha) = 1$: la distribución estacionaria
es $\mathcal{N}(0,1)$, el prior de los modelos de difusión. El proceso OU lleva
cualquier distribución de datos hacia $\mathcal{N}(0,1)$ con velocidad
exponencial. El reverse process (Capítulo~\ref{ch:cap11}) deberá invertir exactamente
este transporte.

Para $\bm{X}_t \in \mathbb{R}^d$ con $d\bm{X}_t = -\bm{X}_t\,dt +
\sqrt{2}\,d\bm{W}_t$, la distribución estacionaria es $\mathcal{N}(\bm{0},
\bm{I})$ componente a componente: el noise schedule de los modelos de difusión más
simples (DDPM) es una discretización de este proceso continuo.
\end{observacion}

\begin{fisicaconexion}{El proceso de Ornstein-Uhlenbeck en física estadística}
El proceso OU fue introducido por Ornstein y Uhlenbeck (1930) para describir
la velocidad de una partícula browniana en un fluido viscoso, corrigiendo la
teoría de Einstein al incluir la inercia. La ecuación de Langevin completa es
$m\dot{v} = -\gamma v + F_{\mathrm{ruido}}(t)$, que en la forma de SDE es
exactamente~\eqref{eq:ou} con $X_t = v_t$, $\alpha = \gamma/m$ y
$\sigma^2 = 2\gamma k_BT/m^2$ (relación fluctuación-disipación). La distribución
estacionaria~\eqref{eq:ou_estacionaria} con $\sigma^2/(2\alpha) = k_BT/m$
es la distribución de Maxwell-Boltzmann para la velocidad: la física
estadística de equilibrio emerge como la distribución estacionaria de la SDE
de Langevin.

La conexión más profunda es con el \textbf{teorema fluctuación-disipación}: el
coeficiente de difusión $\sigma^2$ y la tasa de amortiguamiento $\alpha$ no son
independientes. La relación $\sigma^2 = 2\alpha k_BT/m$ (en unidades apropiadas,
$D = k_BT\mu$ con $\mu$ la movilidad) es necesaria para que la distribución
estacionaria sea el ensemble de Gibbs-Boltzmann $p^*(v) \propto e^{-mv^2/(2k_BT)}$.
Esta es la misma condición de equilibrio que emerge del análisis de la distribución
de Gibbs del SGD del Capítulo~\ref{ch:cap5}: la temperatura del SGD $\eta c/2$
juega el mismo rol que $k_BT/m$ en la ecuación de Langevin física.

En espectroscopía de resonancia magnética nuclear (RMN), el proceso OU describe
la relajación de la magnetización transversal $M_\perp(t) = M_0 e^{-t/T_2}$
bajo la acción del campo fluctuante de los espines vecinos: la tasa $\alpha =
1/T_2$ es el tiempo de relajación espín-espín. Los tiempos de relajación $T_1$
(longitudinal) y $T_2$ (transversal) son directamente los parámetros $\alpha$
de los procesos OU que describen las componentes de la magnetización.
\end{fisicaconexion}

%───────────────────────────────────────────────────────────────────────────────
\section{El generador infinitesimal y la conexión con las EDPs}
\label{sec:generador}
\dificultad{3}{3}
%───────────────────────────────────────────────────────────────────────────────

\subsection{El generador de Markov}
\label{subsec:generador}

\begin{definicion}{Generador infinitesimal}
\label{def:generador}
Para la SDE $dX_t = \mu(X_t)\,dt + \sigma(X_t)\,dW_t$ (autónoma para
simplificar la notación), el \textbf{generador infinitesimal} es el operador
diferencial:
\begin{equation}
  \mathcal{A}f(x)
  = \mu(x)\frac{\partial f}{\partial x}(x)
  + \frac{\sigma(x)^2}{2}\frac{\partial^2 f}{\partial x^2}(x).
  \label{eq:generador}
\end{equation}
\end{definicion}

El nombre se justifica por el siguiente resultado, consecuencia directa de la
fórmula de Itô:

\begin{proposicion}{Caracterización del generador}
\label{prop:generador_def}
Para $f \in C^2$, el proceso $f(X_t) - f(X_0) - \int_0^t \mathcal{A}f(X_s)\,ds$
es una martingala. Equivalentemente:
\[
  \mathcal{A}f(x) = \lim_{t\downarrow 0}
  \frac{\mathbb{E}^x[f(X_t)] - f(x)}{t},
\]
donde $\mathbb{E}^x$ denota la esperanza condicionada a $X_0 = x$.
\end{proposicion}

\begin{proof}
Por la fórmula de Itô: $f(X_t) - f(x) = \int_0^t \mathcal{A}f(X_s)\,ds
+ \int_0^t f'(X_s)\sigma(X_s)\,dW_s$. El último término es una martingala
de Itô con media cero. Tomando la esperanza:
$\mathbb{E}^x[f(X_t)] - f(x) = \int_0^t \mathbb{E}^x[\mathcal{A}f(X_s)]\,ds$.
Para $f, \mu, \sigma$ acotadas, $\mathbb{E}^x[\mathcal{A}f(X_s)] \to \mathcal{A}f(x)$
cuando $s \to 0$, y dividiendo por $t \to 0$ se obtiene el límite. $\square$
\end{proof}

\subsection{La ecuación de Fokker-Planck}
\label{subsec:fokker_planck}

La \textbf{ecuación de Fokker-Planck} describe la evolución de la densidad
de probabilidad $p_t(x) = p(X_t \in dx)/dx$ del proceso $X_t$. Es la ecuación
dual (adjunta $L^2$) de la ecuación de Kolmogorov hacia atrás.

Para un proceso general $dX_t = \mu(X_t,t)\,dt + \sigma(X_t,t)\,dW_t$,
la ecuación de Fokker-Planck es:
\begin{equation}
  \frac{\partial p_t}{\partial t}
  = \mathcal{A}^* p_t
  = -\frac{\partial}{\partial x}(\mu(x,t)\,p_t)
  + \frac{1}{2}\frac{\partial^2}{\partial x^2}(\sigma(x,t)^2\,p_t),
  \label{eq:fp}
\end{equation}
donde $\mathcal{A}^* = -\partial_x \circ \mu + \frac{1}{2}\partial_{xx}
\circ \sigma^2$ es el adjunto $L^2$ del generador $\mathcal{A}$. Esta
ecuación se derivará rigurosamente en el Capítulo~\ref{ch:cap11}.

Para $\mu = 0$ y $\sigma = 1$ (difusión pura):
$\frac{\partial p_t}{\partial t} = \frac{1}{2}\partial_{xx}p_t$. Esta es la
\textbf{ecuación de calor} (con escala temporal $1/2$), que el estudiante
resolvió en física matemática. La densidad de transición del movimiento
browniano es:
\[
  p(X_t = x | X_0 = x_0) = \frac{1}{\sqrt{2\pi t}}e^{-(x-x_0)^2/(2t)},
\]
que es exactamente la función de Green de la ecuación de calor con
fuente puntual en $x_0$: $\partial_t G = \frac{1}{2}\partial_{xx}G$,
$G(x,0;x_0) = \delta(x-x_0)$. La ecuación de calor y el movimiento browniano
son la misma estructura matemática.

\begin{fisicaconexion}{La ecuación de calor, la difusión y la función de Green}
La ecuación de calor $\partial_t T = \kappa\nabla^2 T$ (con $\kappa$ la
difusividad térmica) es uno de los prototipos de EDP parabólica que el físico
resuelve con series de Fourier, separación de variables o transformadas. La
conexión estocástica añade una capa de significado: la solución de la ecuación
de calor con condición inicial $T(\bm{x},0) = T_0(\bm{x})$ puede escribirse
como esperanza sobre el movimiento browniano escalado:
\[
  T(\bm{x},t) = \mathbb{E}^{\bm{x}}[T_0(\bm{x} + \sqrt{2\kappa}\,\bm{W}_t)],
\]
donde el browniano modela la trayectoria aleatoria de un portador de calor.
Este es el caso $V=0$, $g = T_0$, $\mu = 0$, $\sigma = \sqrt{2\kappa}$ de la
fórmula de Feynman-Kac~\eqref{eq:feynman_kac}. La función de Green
$G(\bm{x},t;\bm{x}_0) = (4\pi\kappa t)^{-d/2}e^{-|\bm{x}-\bm{x}_0|^2/(4\kappa t)}$
es la densidad de transición del browniano escalado: temperatura y probabilidad
son la misma función en dos idiomas.

Para la \textbf{ecuación de Schrödinger} en tiempo imaginario $\tau = it$:
$\partial_\tau\psi = -\hat{H}\psi/\hbar$ con $\hat{H} = -(\hbar^2/2m)\nabla^2
+ V(\bm{x})$, la sustitución $\tau \to t$ convierte la ecuación de Schrödinger
en la ecuación de Fokker-Planck~\eqref{eq:fp} (con deriva y potencial). La
densidad de probabilidad cuántica $|\psi|^2$ evoluciona como la densidad de
probabilidad de una difusión. El camino de Feynman es exactamente la
representación de Feynman-Kac~\eqref{eq:feynman_kac} de la función de Green
cuántica: la integral de camino es la integral de Wiener sobre trayectorias
brownianas. Esta dualidad entre mecánica cuántica y procesos estocásticos
(mecánica estocástica de Nelson, 1966) es el fundamento teórico profundo de la
relación entre los modelos de difusión del Capítulo~\ref{ch:cap12} y la física
estadística cuántica.
\end{fisicaconexion}
\label{subsec:kolmogorov_atras}

Sea $u(x,t) = \mathbb{E}^{x,t}[g(X_T)]$ la esperanza condicional de una
función terminal $g$ dado $X_t = x$. La \textbf{ecuación de Kolmogorov hacia
atrás} establece:
\begin{equation}
  \frac{\partial u}{\partial t} + \mathcal{A}u = 0,
  \qquad u(x,T) = g(x).
  \label{eq:kolmogorov_atras}
\end{equation}
Esta ecuación retrocede en tiempo desde la condición terminal $u(x,T) = g(x)$.
La similitud con el backward pass: la ecuación hacia atrás propaga información
del futuro (la función objetivo) al pasado (las condiciones de estado).

\begin{proof}[Derivación de~\eqref{eq:kolmogorov_atras}]
Para $s < T$, por la propiedad de Markov:
$u(X_s,s) = \mathbb{E}[g(X_T)|\mathcal{F}_s]$.
El proceso $u(X_s,s)$ es una martingala (esperanza condicional de $g(X_T)$).
Por la fórmula de Itô~\eqref{eq:ito_formula}:
\[
  du(X_s,s) = \left(\partial_t u + \mathcal{A}u\right)ds
  + \sigma(X_s,s)\partial_x u\,dW_s.
\]
Para que $u(X_s,s)$ sea martingala, el coeficiente de $ds$ debe ser cero:
$\partial_t u + \mathcal{A}u = 0$. La condición terminal es $u(x,T) = g(x)$
por definición. $\square$
\end{proof}

\subsection{La fórmula de Feynman-Kac}
\label{subsec:feynman_kac}

La fórmula de Feynman-Kac generaliza la conexión entre EDPs y esperanzas
estocásticas al caso con potencial. Para la EDP:
\begin{equation}
  \frac{\partial u}{\partial t} + \mathcal{A}u - V(x,t)u = 0,
  \qquad u(x,T) = g(x),
  \label{eq:feynman_kac_pde}
\end{equation}
con $V \geq 0$ un potencial, la solución es:

\begin{teorema}{Fórmula de Feynman-Kac}
\label{thm:feynman_kac}
Bajo condiciones de regularidad apropiadas sobre $\mu$, $\sigma$, $V$ y $g$,
la solución de la EDP~\eqref{eq:feynman_kac_pde} es:
\begin{equation}
  u(x,t)
  = \mathbb{E}^{x,t}\!\left[
    \exp\!\left(-\int_t^T V(X_s,s)\,ds\right) g(X_T)
  \right].
  \label{eq:feynman_kac}
\end{equation}
\end{teorema}

\begin{proof}
Definir el proceso adjunto con factor de descuento:
$M_s = e^{-\int_t^s V(X_r,r)\,dr}\,u(X_s,s)$ para $s \in [t,T]$.
Se demostrará que $M_s$ es martingala, lo que implica
$\mathbb{E}^{x,t}[M_T] = M_t = u(x,t)$.

Aplicar la fórmula de Itô al producto:
\[
  dM_s = e^{-\int_t^s V\,dr}\left[
    (-V(X_s,s)u + \partial_t u + \mathcal{A}u)\,ds
    + \sigma\partial_x u\,dW_s
  \right].
\]
Usando~\eqref{eq:feynman_kac_pde}: $-Vu + \partial_t u + \mathcal{A}u = 0$.
Por tanto $dM_s = e^{-\int_t^s V\,dr}\sigma(X_s,s)\partial_x u(X_s,s)\,dW_s$:
el coeficiente de $ds$ es cero y $M_s$ es martingala.

Tomando la esperanza condicional con $M_t = u(x,t)$ y usando la condición
terminal $u(X_T,T) = g(X_T)$:
$u(x,t) = \mathbb{E}^{x,t}[M_T]
= \mathbb{E}^{x,t}[e^{-\int_t^T V\,dr}\,g(X_T)]$. $\square$
\end{proof}

\begin{observacion}{Feynman-Kac y la función de score}
En el Capítulo~\ref{ch:cap12}, la función de score $\nabla_{\bm{x}}\log p_t(\bm{x})$
de los modelos de difusión satisface una EDP parabólica de la familia
\eqref{eq:feynman_kac_pde}. La fórmula de Feynman-Kac~\eqref{eq:feynman_kac}
expresa el score como una esperanza sobre el proceso estocástico forward,
lo que proporciona la base teórica del \textit{score matching} para
entrenar modelos de difusión: en lugar de resolver la EDP directamente,
se estima la esperanza estocástica por Monte Carlo.
\end{observacion}

%───────────────────────────────────────────────────────────────────────────────
\section{Resumen del Capítulo 10}
\label{sec:resumen10}
%───────────────────────────────────────────────────────────────────────────────

Este capítulo construyó el cálculo de Itô sobre tres pilares. El primero es la
variación cuadrática del movimiento browniano: $[W,W]_T = T$ c.s., resultado
demostrado con convergencia casi segura vía el segundo lema de Borel-Cantelli.
Este resultado es la razón precisa de $(dW)^2 = dt$ y la fuente de todas las
diferencias entre el cálculo ordinario y el estocástico. El segundo pilar es
la integral de Itô, construida como límite $L^2$ de sumas de Riemann con punto
izquierdo, con la isometría de Itô como herramienta central. El tercer pilar es
la fórmula de Itô, demostrada por expansión de Taylor reteniedo el término de
segundo orden en $dX$, que produce la corrección $\frac{\sigma^2}{2}\partial_{xx}f$
inexistente en el cálculo ordinario.

El hilo hacia la física matemática del estudiante es la ecuación de calor:
la difusión pura $dX_t = dW_t$ tiene como densidad de transición exactamente
la función de Green de la ecuación de calor $\partial_t p = \frac{1}{2}\Delta p$.
Las EDPs parabólicas y las SDEs son la misma estructura matemática: las
soluciones de EDPs son esperanzas de funcionales estocásticos (Feynman-Kac)
y las densidades de transición de SDEs son soluciones de EDPs (Fokker-Planck).

Las dos conexiones hacia adelante son precisas. El \textbf{Capítulo~\ref{ch:cap11}}
derivará rigurosamente la ecuación de Fokker-Planck a partir de la fórmula
de Itô, estudiará su distribución estacionaria, y demostrará el teorema de
Anderson para el time reversal: la SDE cuya densidad de transición va hacia
atrás en tiempo, con la función de score como término de deriva adicional.
La demostración del teorema de Anderson usa exactamente la fórmula de Itô
y las propiedades de martingala desarrolladas aquí. El \textbf{Capítulo~\ref{ch:cap12}}
usará el proceso OU como forward process y el teorema de Anderson como
reverse process para construir los modelos de difusión: el score network
entrenado es la función de score de Feynman-Kac estimada por matching.

%───────────────────────────────────────────────────────────────────────────────
\label{sec:estrella10}
\begin{estrella}{Simulación estocástica en reactores de investigación}

\textbf{La SDE de Langevin para la cinética puntual con ruido.}
La ecuación de cinética puntual del reactor (Capítulo~\ref{ch:cap5}) asume
que la reactividad $\rho(t)$ es una función determinista del tiempo. En la
práctica, el flujo neutrónico en el IAN-R1 fluctúa estadísticamente incluso
en condiciones nominales de operación, por dos fuentes de ruido: (1) la
naturaleza discreta de los eventos de fisión (ruido de Poisson de las fuentes
neutrónicas) y (2) las fluctuaciones mecánicas de las barras de control
(vibraciones). El modelo estocástico del flujo neutrónico en estado estacionario
perturbado es:
\[
  d\Phi_t = \bigl[-\alpha(\Phi_t - \Phi_0)\bigr]\,dt
  + \sigma_\Phi\,dW_t,
\]
que es exactamente el proceso OU~\eqref{eq:ou} con $X_t = \Phi_t$, tasa de
retorno $\alpha$ determinada por el tiempo de vida de los neutrones prompts, y
$\Phi_0$ el flujo de referencia. La solución explícita~\eqref{eq:ou_sol}
predice que las fluctuaciones de flujo tienen distribución gaussiana con
varianza $\sigma_\Phi^2/(2\alpha)$ en estado estacionario.

\textbf{Medición experimental de $\alpha$ por análisis de ruido.}
La técnica de análisis de ruido en reactores (Noise Analysis) mide la función
de autocorrelación del flujo:
\[
  R_\Phi(\tau) = \mathbb{E}[\Phi_t\Phi_{t+\tau}]
  = \Phi_0^2 + \frac{\sigma_\Phi^2}{2\alpha}e^{-\alpha|\tau|},
\]
que se obtiene directamente de la solución~\eqref{eq:ou_sol} y las propiedades
del proceso OU. Midiendo $R_\Phi(\tau)$ experimentalmente (correlación de
señales de cámara de ionización) y ajustando la exponencial decayente, se
extrae $\alpha$ sin perturbar el reactor. Este método es equivalente al
análisis de Rossi-$\alpha$: ambos miden la tasa de decaimiento $\alpha$ del
modo prompt del reactor. La isometría de Itô (Teorema~\ref{thm:isometria})
garantiza que la varianza de la estimación de $\alpha$ puede acotarse en
función del tiempo de observación $T$ y la varianza del proceso.

\textbf{El proceso OU multidimensional para barras acopladas.}
El IAN-R1 tiene varias barras de control cuyas posiciones fluctúan por
vibraciones. El modelo completo es:
\[
  d\bm{X}_t = -\bm{A}\bm{X}_t\,dt + \bm{\Sigma}\,d\bm{W}_t,
\]
donde $\bm{X}_t \in \mathbb{R}^d$ son las posiciones de las barras y $\bm{A}$
es la matriz de rigidez mecánica. La fórmula de Itô multidimensional~\eqref{eq:ito_multi}
con $f = \Phi(\bm{X})$ el mapa barra$\to$flujo (sustituto MLP del Capítulo~\ref{ch:cap8})
da directamente la SDE para el flujo inducido:
\[
  d\Phi_t = \nabla_{\bm{X}}\Phi\cdot(-\bm{A}\bm{X}_t)\,dt
  + \frac{1}{2}\mathrm{tr}(\bm{\Sigma}\bm{\Sigma}^\top\nabla^2_{\bm{X}}\Phi)\,dt
  + (\nabla_{\bm{X}}\Phi)^\top\bm{\Sigma}\,d\bm{W}_t.
\]
El término de corrección de Itô $\frac{1}{2}\mathrm{tr}(\bm{\Sigma}\bm{\Sigma}^\top
\nabla^2_{\bm{X}}\Phi)$ es la corrección de segundo orden que aparece por la
no linealidad del mapa sustituto: el flujo medio no es simplemente
$\Phi(\mathbb{E}[\bm{X}_t])$, sino que hay una corrección de Jensen debida a
la curvatura del sustituto, exactamente el término $\sigma^2\partial_{xx}f/2$
de la fórmula de Itô.
\end{estrella}

%───────────────────────────────────────────────────────────────────────────────
\section*{Ejercicios computacionales}
\addcontentsline{toc}{section}{Ejercicios computacionales}
%───────────────────────────────────────────────────────────────────────────────

\begin{ejerciciocomp}{Verificación de la variación cuadrática y la isometría de Itô}
\textbf{Objetivo:} verificar numéricamente el Teorema~\ref{thm:vc} y la
isometría de Itô~\eqref{eq:isometria} por simulación de trayectorias brownianas.

\medskip\noindent\textbf{Algoritmo:}
\begin{enumerate}
  \item Simular $N = 10^3$ trayectorias del browniano estándar en $[0,1]$
        con $n \in \{10, 100, 1000\}$ subintervalos (incrementos
        $\Delta W_k \sim \mathcal{N}(0, 1/n)$ i.i.d.).
  \item Para cada $n$, calcular $V_1^{(2)}(\Pi_n) = \sum_k(\Delta W_k)^2$
        por trayectoria. Verificar que $\mathbb{E}[V_1^{(2)}] \approx 1$ y
        $\mathrm{Var}[V_1^{(2)}] \approx 2/n$ (consistente con la demostración
        del Teorema~\ref{thm:vc}).
  \item Calcular la integral de Itô $I = \int_0^1 W_t\,dW_t = \sum_k W_{t_k}\Delta W_k$
        por simulación. Verificar la isometría: $\mathbb{E}[I^2] \approx
        \int_0^1 \mathbb{E}[W_t^2]\,dt = \int_0^1 t\,dt = 1/2$.
  \item Verificar la fórmula de Itô: $\int_0^1 W_t\,dW_t = (W_1^2 - 1)/2$
        (caso $f(x) = x^2/2$). Comparar el resultado de la suma de Riemann
        estocástica con la fórmula exacta para cada trayectoria.
\end{enumerate}
\medskip\noindent\textbf{Verificación:} el histograma de $I = (W_1^2-1)/2$
debe ser una distribución $(\chi^2(1) - 1)/2$.
\end{ejerciciocomp}

\begin{ejerciciocomp}{Simulación del proceso OU y estimación de parámetros}
\textbf{Objetivo:} simular el proceso OU, verificar la distribución estacionaria
y estimar los parámetros $\alpha$ y $\sigma$ a partir de trayectorias ruidosas
(análogo al análisis de ruido en reactores).

\medskip\noindent\textbf{Algoritmo:}
\begin{enumerate}
  \item Simular $N = 500$ trayectorias del OU con $\alpha = 2$, $\sigma = 1$,
        $X_0 \sim \mathcal{N}(0,1)$ y paso $\Delta t = 0.01$ hasta $T = 10$
        usando Euler-Maruyama: $X_{k+1} = X_k - \alpha X_k\Delta t +
        \sigma\sqrt{\Delta t}\,\xi_k$.
  \item Verificar que el histograma de $X_T$ (con $T$ grande) aproxima
        $\mathcal{N}(0, \sigma^2/(2\alpha)) = \mathcal{N}(0, 0.25)$.
  \item Estimar $\alpha$ desde la función de autocorrelación empírica
        $\hat{R}(\tau) = \frac{1}{T}\int_0^{T-\tau}X_t X_{t+\tau}\,dt$:
        ajustar $\hat{R}(\tau) = \hat{\sigma}^2_\infty e^{-\hat{\alpha}|\tau|}$
        por mínimos cuadrados sobre $\tau \in [0, 2]$.
  \item Comparar los parámetros estimados $(\hat{\alpha}, \hat{\sigma})$ con
        los verdaderos en función del número de trayectorias $N$ y el tiempo
        de observación $T$. ¿Cuál es más importante para la estimación precisa?
\end{enumerate}
\end{ejerciciocomp}

%───────────────────────────────────────────────────────────────────────────────
\begin{notasbib}
La construcción rigurosa del movimiento browniano fue realizada por Wiener
(1923)~\cite{wiener1923differential}, quien construyó la medida de Wiener sobre
el espacio $C([0,T])$. El teorema de Donsker (Sección~\ref{subsec:caminata})
es la versión funcional del teorema central del límite; la referencia estándar
es Billingsley (1999)~\cite{billingsley1999convergence}. El teorema de Kolmogorov-Chentsov
para continuidad de trayectorias es estándar en cualquier texto de procesos
estocásticos; véase Karatzas y Shreve (1991)~\cite{karatzas1991brownian}, la
referencia canónica del campo.

La integral de Itô fue introducida por Kiyosi Itô (1944)~\cite{ito1944stochastic}.
La isometría de Itô (Teorema~\ref{thm:isometria}) y la extensión por densidad
son la construcción estándar; véase Øksendal (2003)~\cite{oksendal2003stochastic}
para una presentación accesible. La integral de Stratonovich (Sección~\ref{subsec:stratonovich})
fue introducida por Stratonovich (1966)~\cite{stratonovich1966new}; la conversión
entre las dos formulaciones~\eqref{eq:ito_strat} es estándar.

La fórmula de Itô (Teorema~\ref{thm:ito_formula}) fue demostrada por Itô
(1951)~\cite{ito1951formula}. El movimiento browniano geométrico (Ejercicio
Resuelto~\ref{ej:browniano_geometrico}) es el modelo de Black y Scholes (1973)
para precios de opciones; la fórmula de Black-Scholes es una aplicación directa
de la fórmula de Itô a la EDP de Kolmogorov hacia atrás~\eqref{eq:kolmogorov_atras}.

El proceso de Ornstein-Uhlenbeck fue introducido por Ornstein y Uhlenbeck
(1930)~\cite{ornstein1930brownian} para modelar la velocidad browniana con
inercia. Su solución explícita~\eqref{eq:ou_sol} y la distribución
estacionaria~\eqref{eq:ou_estacionaria} son clásicas. El teorema de existencia
y unicidad (Teorema~\ref{thm:ito_doob}) es debido a Itô y Doob; la prueba por
iteración de Picard sigue a Øksendal (2003).

La ecuación de Fokker-Planck (enunciada en la Sección~\ref{subsec:fokker_planck}
y a derivar en el Capítulo~\ref{ch:cap11}) fue introducida por Fokker (1914) y
Planck (1917) en el contexto de la difusión browniana. La ecuación de Kolmogorov
hacia atrás~\eqref{eq:kolmogorov_atras} y la fórmula de Feynman-Kac
(Teorema~\ref{thm:feynman_kac}) son debidas a Kolmogorov (1931) y a Feynman
(1948) y Kac (1949)~\cite{kac1949distributions} respectivamente. La conexión
entre la integral de camino de Feynman y la medida de Wiener, mencionada en la
\textsf{Conexión física} de la Sección~\ref{subsec:fokker_planck}, fue
desarrollada por Nelson (1966)~\cite{nelson1966derivation}.

Para el análisis de ruido en reactores nucleares (Sección~\ref{sec:estrella10}),
la referencia estándar es Williams (1974)~\cite{williams1974random}; el análisis
de Rossi-$\alpha$ y los métodos de correlación para medir parámetros cinéticos
se revisan en Pázsit y Pál (2008)~\cite{pazsit2008neutron}.
\end{notasbib}

%───────────────────────────────────────────────────────────────────────────────
\section*{Ejercicios resueltos}
\addcontentsline{toc}{section}{Ejercicios resueltos}
%───────────────────────────────────────────────────────────────────────────────

\begin{ejercicio}
\label{ej:martingalas}
\textbf{(Verificación de martingalas del browniano).}

Verificar directamente desde la definición (sin usar la fórmula de Itô) que
$M_t^{(2)} = W_t^2 - t$ es martingala.

\medskip\noindent\textit{Solución.}
Para $s < t$, usando los incrementos independientes del browniano:
\begin{align*}
  \mathbb{E}[W_t^2|\mathcal{F}_s]
  &= \mathbb{E}[(W_s + (W_t-W_s))^2|\mathcal{F}_s] \\
  &= \mathbb{E}[W_s^2 + 2W_s(W_t-W_s) + (W_t-W_s)^2|\mathcal{F}_s] \\
  &= W_s^2 + 2W_s\,\mathbb{E}[W_t-W_s|\mathcal{F}_s]
    + \mathbb{E}[(W_t-W_s)^2|\mathcal{F}_s].
\end{align*}
El incremento $W_t-W_s$ es independiente de $\mathcal{F}_s$:
$\mathbb{E}[W_t-W_s|\mathcal{F}_s] = \mathbb{E}[W_t-W_s] = 0$ y
$\mathbb{E}[(W_t-W_s)^2|\mathcal{F}_s] = \mathrm{Var}[W_t-W_s] = t-s$.
Por tanto:
\[
  \mathbb{E}[M_t^{(2)}|\mathcal{F}_s]
  = \mathbb{E}[W_t^2-t|\mathcal{F}_s]
  = W_s^2 + (t-s) - t = W_s^2 - s = M_s^{(2)}.
\]
$M_t^{(2)}$ es martingala. $\square$
\end{ejercicio}

\begin{ejercicio}
\label{ej:browniano_geometrico}
\textbf{(Movimiento browniano geométrico: solución explícita y paradoja de Itô).}

Sea $dS_t = \mu S_t\,dt + \sigma S_t\,dW_t$ con $S_0 = s_0 > 0$.

\textbf{(a)} Encontrar la solución explícita aplicando la fórmula de Itô a
$f(S) = \log S$.

\textbf{(b)} Verificar que la solución preserva $S_t > 0$ c.s.

\textbf{(c)} La "paradoja": la solución tiene $\mathbb{E}[S_t] = s_0 e^{\mu t}$,
que crece con tasa $\mu$. Pero el logaritmo de la media geométrica satisface
$\mathbb{E}[\log S_t] = \log s_0 + (\mu - \sigma^2/2)t$, que crece con tasa
$\mu - \sigma^2/2 < \mu$. ¿Por qué la tasa "efectiva" es menor que $\mu$?

\medskip\noindent\textit{Solución.}
\textbf{(a)} Aplicar la fórmula de Itô a $f(S) = \log S$, con $\partial_S f = 1/S$,
$\partial_{SS}f = -1/S^2$, y $dS_t = \mu S_t\,dt + \sigma S_t\,dW_t$:
\[
  d\log S_t
  = \frac{1}{S_t}\,dS_t - \frac{1}{2S_t^2}(dS_t)^2
  = (\mu - \tfrac{\sigma^2}{2})\,dt + \sigma\,dW_t.
\]
Integrando: $\log S_t = \log s_0 + (\mu-\sigma^2/2)t + \sigma W_t$, es decir:
$S_t = s_0\exp((\mu-\sigma^2/2)t + \sigma W_t)$.

\textbf{(b)} Como $S_t = s_0\exp(\cdot)$ con $s_0 > 0$, se tiene $S_t > 0$ c.s.

\textbf{(c)} La corrección $-\sigma^2/2$ es la convexidad de la función exponencial
vista desde el teorema de Jensen: $\mathbb{E}[e^X] \geq e^{\mathbb{E}[X]}$ para
cualquier variable aleatoria $X$. La media aritmética $\mathbb{E}[S_t]$ crece
más rápido que la media geométrica $e^{\mathbb{E}[\log S_t]}$ por la varianza
del browniano. El término $-\sigma^2/2$ en la tasa del logaritmo es el precio
estocástico de la convexidad: exactamente la corrección de Itô. $\square$
\end{ejercicio}

\begin{ejercicio}
\label{ej:ou_distribucion}
\textbf{(Distribución transitoria del proceso OU y convergencia a la estacionaria).}

Para el proceso OU $dX_t = -X_t\,dt + \sqrt{2}\,dW_t$ con $X_0 \sim
\mathcal{N}(\mu_0,\sigma_0^2)$:

\textbf{(a)} Calcular $\mathbb{E}[X_t]$ y $\mathrm{Var}[X_t]$ para todo $t \geq 0$.

\textbf{(b)} Mostrar que $X_t \sim \mathcal{N}(\mu_0 e^{-t}, \sigma_0^2 e^{-2t}
+ 1 - e^{-2t})$ para todo $t$.

\textbf{(c)} ¿Para qué condición inicial $X_0$ la distribución de $X_t$ es
estacionaria (idéntica para todo $t$)?

\medskip\noindent\textit{Solución.}
\textbf{(a)} De~\eqref{eq:ou_sol} con $\alpha = 1$, $\sigma = \sqrt{2}$:
$X_t = X_0 e^{-t} + \sqrt{2}\int_0^t e^{-(t-s)}\,dW_s$.
$\mathbb{E}[X_t] = \mu_0 e^{-t}$.
$\mathrm{Var}[X_t] = \sigma_0^2 e^{-2t} + 2\int_0^t e^{-2(t-s)}\,ds
= \sigma_0^2 e^{-2t} + 1 - e^{-2t}$.

\textbf{(b)} $X_t$ es suma de una gaussiana ($X_0 e^{-t}$) y una integral
estocástica de proceso determinista (gaussiana), ambas independientes:
su distribución es gaussiana con la media y varianza calculadas.

\textbf{(c)} La distribución es estacionaria si no depende de $t$:
$\mu_0 e^{-t} = \mu_0$ requiere $\mu_0 = 0$; $\sigma_0^2 e^{-2t}+1-e^{-2t}
= \sigma_0^2$ requiere $\sigma_0^2 = 1$. La condición inicial estacionaria
es $X_0 \sim \mathcal{N}(0,1)$. $\square$
\end{ejercicio}

%───────────────────────────────────────────────────────────────────────────────
\section*{Ejercicios propuestos}
\addcontentsline{toc}{section}{Ejercicios propuestos}
%───────────────────────────────────────────────────────────────────────────────

\begin{ejercicio}
\textbf{(Variación $p$-ésima del movimiento browniano).}
Sea $V_T^{(p)}(\Pi) = \sum_i |\Delta W_i|^p$ la variación $p$-ésima.
\begin{enumerate}[label=(\alph*)]
  \item Para $p = 1$ (variación total): mostrar que $\mathbb{E}[V_T^{(1)}]
        = \sqrt{2/\pi}\sum_i\sqrt{\Delta t_i} \sim \sqrt{n}$ para partición
        uniforme con $n$ subintervalos. ¿Qué ocurre cuando $n \to \infty$?
  \item Para $p = 2$: verificar el Teorema~\ref{thm:vc} calculando
        $\mathrm{Var}[V_T^{(2)}]$ explícitamente usando $\mathbb{E}[Z^4] = 3$
        para $Z \sim \mathcal{N}(0,1)$.
  \item Para $p > 2$: mostrar que $V_T^{(p)} \to 0$ en probabilidad cuando
        $\|\Pi\| \to 0$. ¿Qué ocurre para $p < 2$?
\end{enumerate}
\end{ejercicio}

\begin{ejercicio}
\textbf{(Isometría de Itô e independencia del punto de evaluación).}
Definir la integral "de punto derecho"
$\int_0^T f(t)\,dW_t^{\mathrm{right}} = \lim_{\|\Pi\|\to 0}\sum_i f(t_{i+1})
\Delta W_i$ (análoga a la de Itô pero usando el final del intervalo).
\begin{enumerate}[label=(\alph*)]
  \item ¿Satisface la integral de punto derecho la isometría de Itô?
  \item Calcular $\int_0^T W_t\,dW_t$ con la integral de Itô y con la de
        punto derecho. ¿Difieren?
  \item Mostrar que la diferencia entre las dos integrales es exactamente el
        término de corrección de Stratonovich~\eqref{eq:ito_strat}.
\end{enumerate}
\end{ejercicio}

\begin{ejercicio}
\textbf{(Fórmula de Itô para el proceso de CIR).}
El proceso de Cox-Ingersoll-Ross (CIR): $dX_t = \kappa(\theta-X_t)\,dt +
\sigma\sqrt{X_t}\,dW_t$ con $X_0 = x_0 > 0$.
\begin{enumerate}[label=(\alph*)]
  \item Aplicar la fórmula de Itô a $f(X) = X^2$ para obtener la evolución
        de $\mathbb{E}[X_t^2]$, y resolver la ODE resultante.
  \item Para $\kappa = 1$, $\theta = 1$, $\sigma = 0.5$, calcular
        $\mathrm{Var}[X_t]$ para $t \in [0,5]$ y verificar que $X_t \geq 0$
        c.s.\ si $2\kappa\theta \geq \sigma^2$ (condición de Feller).
  \item Comparar la distribución estacionaria del CIR con la del proceso OU:
        ¿cuál es la diferencia cualitativa más importante?
\end{enumerate}
\end{ejercicio}

\begin{ejercicio}
\textbf{(Ecuación de Kolmogorov y la ecuación de calor).}
Sea $X_t$ el movimiento browniano estándar y $u(x,t) = \mathbb{E}^{x,0}
[\sin(X_T)]$ para $T$ fijo.
\begin{enumerate}[label=(\alph*)]
  \item Escribir la ecuación de Kolmogorov hacia atrás para $u$.
  \item Resolver la ecuación (con condición terminal $u(x,T) = \sin x$) usando
        la transformada de Fourier.
  \item Verificar la solución obtenida con la fórmula de Feynman-Kac
        ($V = 0$): calcular $\mathbb{E}[\sin(X_T)|X_0 = x]$ directamente usando
        la función generatriz del browniano.
  \item Interpretar el factor de amortiguamiento en la solución: ¿por qué
        $u(x,0)$ tiene menor amplitud que $\sin x$?
\end{enumerate}
\end{ejercicio}

\begin{ejercicio}
\textbf{(Simulación de SDEs: el método de Euler-Maruyama).}
El método de Euler-Maruyama discretiza $dX_t = \mu(X_t)\,dt + \sigma(X_t)\,dW_t$
con paso $\Delta t$:
$X_{k+1} = X_k + \mu(X_k)\Delta t + \sigma(X_k)\sqrt{\Delta t}\,\xi_k$,
$\xi_k \sim \mathcal{N}(0,1)$ i.i.d.
\begin{enumerate}[label=(\alph*)]
  \item Implementar Euler-Maruyama para el proceso OU del
        Ejercicio~\ref{ej:ou_distribucion} con $\Delta t \in \{0.1, 0.01, 0.001\}$.
  \item Comparar la distribución empírica de $10^4$ trayectorias en $t = 1$
        con la distribución analítica del Ejercicio Resuelto~\ref{ej:ou_distribucion}.
        ¿Cuál es el error como función de $\Delta t$?
  \item Verificar que la distribución estacionaria empírica (histograma de
        $X_t$ para $t$ grande) converge a $\mathcal{N}(0,1)$.
  \item El método de Milstein añade una corrección de segundo orden:
        $X_{k+1} = X_k + \mu(X_k)\Delta t + \sigma(X_k)\sqrt{\Delta t}\,\xi_k
        + \frac{1}{2}\sigma(X_k)\sigma'(X_k)(\xi_k^2-1)\Delta t$.
        Comparar la precisión de Euler-Maruyama y Milstein para el browniano
        geométrico (Ejercicio Resuelto~\ref{ej:browniano_geometrico}).
\end{enumerate}
\end{ejercicio}

\printbibliography[heading=subbibliography]